\documentclass[11pt,a4paper]{article}

% ============================================================
% Cuarto Informe de Avance -- ATLAS II -- Componente Solar
% Versión Consolidada y Completa
% ============================================================

\usepackage[utf8]{inputenc}
\usepackage[T1]{fontenc}
\usepackage[spanish,es-nodecimaldot]{babel}
\usepackage{geometry}
\geometry{margin=2.5cm}
\usepackage{amsmath,amssymb}
\usepackage{graphicx}
\usepackage{booktabs}
\usepackage{longtable}
\usepackage{array}
\usepackage{enumitem}
\usepackage{xcolor}
\usepackage{hyperref}
\usepackage{caption}
\usepackage{float}

\hypersetup{
    colorlinks=true,
    linkcolor=blue!55!black,
    citecolor=blue!55!black,
    urlcolor=blue!55!black
}

\setlength{\parindent}{0pt}
\setlength{\parskip}{0.65em}

% ---------- Marcadores para completar cuando cierre el procesamiento ----------
\newcommand{\pendiente}[1]{\textcolor{red!70!black}{\textbf{[PENDIENTE: #1]}}}

% ---------- Figuras que no rompen la compilación si todavía no existen ----------
\newcommand{\figuraopcional}[4][0.95\textwidth]{%
\begin{figure}[H]
    \centering
    \IfFileExists{#2}{%
        \includegraphics[width=#1]{#2}%
    }{%
        \IfFileExists{reportes/#2}{%
            \includegraphics[width=#1]{reportes/#2}%
        }{%
            \IfFileExists{../#2}{%
                \includegraphics[width=#1]{../#2}%
            }{%
                \fbox{\parbox[c][5cm][c]{0.88\textwidth}{\centering
                Figura disponible en\\
                \texttt{\detokenize{#2}}}}%
            }%
        }%
    }
    \caption{#3}
    \label{#4}
\end{figure}
}

\title{\textbf{Convenio ATLAS II -- Provincia de Córdoba}\\
Componente Solar\\[1.2em]
\Large \textbf{Cuarto Informe de Avance}\\[0.4em]
\large Atlas Solar Territorial, Potencial Fotovoltaico y Ganancia por Seguimiento\\
\normalsize Estadísticas multianuales del período 2020--2025 sobre grilla territorial nominal de 2 km}
\author{Facultad de Ciencias Exactas, Físicas y Naturales (FCEFyN)\\
Laboratorio de Hidrometeorología}
\date{Octubre de 2026}

\begin{document}

\maketitle
\thispagestyle{empty}
\newpage

% ============================================================
\section*{Resumen general}
% ============================================================

El presente informe documenta la construcción del producto territorial del Atlas Solar de la Provincia de Córdoba a partir de la arquitectura satelital desarrollada y validada en las etapas anteriores del convenio. El objetivo de esta fase es transformar el modelo puntual satélite--superficie en campos espaciales continuos del recurso solar y, a partir de ellos, obtener indicadores derivados de potencial fotovoltaico bajo configuraciones tecnológicas estandarizadas.

La base temporal utilizada comprende seis años completos, correspondientes al período 2020--2025. Sobre la grilla territorial definida para el proyecto se reconstruyeron series de irradiancia global horizontal (GHI) mediante el modelo híbrido BD-JPT + Histogram Gradient Boosting (HGB), formulado en el dominio del índice de claridad. La inferencia territorial incorpora la regionalización solar de la red de superficie y una asignación espacial continua mediante vecinos próximos ponderados por distancia, evitando discontinuidades artificiales entre regiones operativas.

A partir de las series de GHI se calcularon estadísticas multianuales del período 2020--2025 y productos derivados de irradiancia e irradiación. La irradiación directa normal (DNI) y la irradiación difusa horizontal (DHI) se estimaron mediante el modelo de descomposición de Erbs y, posteriormente, las componentes radiativas fueron trasladadas al plano del generador mediante el modelo anisotrópico de Perez. Se evaluaron dos configuraciones fotovoltaicas de referencia: un sistema fijo orientado al norte con inclinación de $30^\circ$ y un seguidor horizontal de un eje con eje norte--sur y backtracking.

La conversión desde irradiación en plano del generador hacia producción eléctrica potencial se realizó mediante un esquema normalizado a 1~kWp, utilizando un modelo térmico NOCT y la formulación PVWatts para la potencia fotovoltaica. Los productos resultantes incluyen GHI, DNI, DHI, horas sol pico, irradiación en plano del generador, rendimiento final, factor de capacidad y ganancia relativa del seguimiento.

Por la extensión temporal disponible, los resultados se presentan como \textbf{estadísticas multianuales observadas/modeladas del período 2020--2025}. No se interpretan como una normal climatológica de largo plazo. Los valores numéricos consolidados y las tablas regionales reflejan la base de datos de alta resolución a 2~km.

\newpage
\tableofcontents
\newpage

% ===============================================================
% 1. INTRODUCCIÓN
% ===============================================================

\section{Introducción}

\subsection{Objetivo}

El Convenio ATLAS II, Componente Solar, tiene como objetivo generar una representación espacial y temporal consistente del recurso solar de la Provincia de Córdoba mediante la integración de observaciones satelitales geoestacionarias, mediciones terrestres calibradas y modelos físicos y estadísticos de irradiancia.

El presente Cuarto Informe aborda la etapa territorial del proyecto. A diferencia de los informes anteriores, centrados en la preparación de la base observacional, la calibración de la red terrestre y la construcción del modelo satelital, esta fase se concentra en la aplicación sistemática del modelo sobre toda la provincia, la construcción de mapas multianuales del recurso solar y la estimación de indicadores de potencial fotovoltaico.

Los objetivos específicos son:

\begin{enumerate}[label=\arabic*)]
    \item aplicar el modelo satelital híbrido sobre la grilla territorial del Atlas para el período completo 2020--2025;
    \item construir campos espaciales continuos de GHI y sus agregaciones temporales;
    \item estimar DNI y DHI como productos derivados de GHI;
    \item calcular irradiación sobre plano inclinado para configuraciones fotovoltaicas fijas y con seguimiento;
    \item transformar el recurso radiativo en indicadores normalizados de producción eléctrica potencial;
    \item cuantificar la ganancia relativa asociada al seguimiento de un eje;
    \item analizar la distribución espacial y regional de los productos obtenidos;
    \item consolidar los productos geoespaciales finales del Atlas.
\end{enumerate}

\subsection{Antecedentes del proyecto}

En el Primer Informe de Avance se estableció el dominio espacial, la grilla objetivo, la estructura reproducible del proyecto y el marco físico-empírico basado en observaciones GOES y modelos tipo JPT/BD-JPT. En el Segundo Informe se documentó la consolidación de la red terrestre para calibración, incorporando el control de calidad, la simulación de cielo despejado mediante REST2 \cite{gueymard2008} y la calibración de los sensores Davis. En el Tercer Informe se documentó la integración entre ambas fuentes, la reformulación del problema en el dominio del índice de claridad ($K_t$) y la selección de una arquitectura híbrida BD-JPT + HGB como núcleo predictivo del Atlas.

Durante la consolidación previa a la inferencia territorial se revisaron nuevamente las convenciones espaciales y temporales de la cadena de procesamiento y se regeneraron los productos intermedios necesarios para garantizar una base homogénea. Por lo tanto, las configuraciones y métricas numéricas presentadas en este informe corresponden a la cadena consolidada utilizada para construir el Atlas territorial; cuando difieran de valores intermedios documentados en informes previos, los resultados del presente informe deben considerarse la versión operativa actualizada.

El objetivo de este informe no es repetir el desarrollo de la calibración satelital, sino documentar cómo el modelo validado se transforma en un producto cartográfico y energético provincial.

\subsection{Alcance temporal}

La ventana de análisis comprende seis años completos:

\[
2020,\;2021,\;2022,\;2023,\;2024,\;2025.
\]

Los estadísticos agregados resumen el comportamiento medio observado y modelado durante los seis años analizados.

\begin{table}[htbp]
\centering
\small
\caption{Métricas provinciales consolidadas del período 2020--2025.}
\label{tab:metricas_2020_2025}
\begin{tabular}{lrrrr}
\toprule
Variable & Media provincial & Mínimo & Máximo & Unidad \\
\midrule
GHI & 1889.3 & 1081.9 & 2054.7 & kWh/m$^2$/año \\
DNI & 1985.2 & 396.6 & 2279.2 & kWh/m$^2$/año \\
DHI & 604.9 & 564.5 & 910.2 & kWh/m$^2$/año \\
HSP & 5.17 & 2.96 & 5.63 & HSP/día \\
POA fijo & 2016.6 & 1042.3 & 2206.0 & kWh/m$^2$/año \\
POA tracker & 2461.6 & 1201.5 & 2701.6 & kWh/m$^2$/año \\
$Y_f$ fijo & 1644.7 & 870.5 & 1792.6 & kWh/kWp/año \\
$Y_f$ tracker & 1982.2 & 995.6 & 2167.5 & kWh/kWp/año \\
CF fijo & 18.77 & 9.94 & 20.46 & \% \\
CF tracker & 22.63 & 11.37 & 24.74 & \% \\
Ganancia tracker & 20.51 & 14.38 & 23.36 & \% \\
CV interanual & 2.14 & 0.47 & 16.99 & \% \\
\bottomrule
\end{tabular}
\end{table}


% ===============================================================
% 2. CONSTRUCCIÓN METODOLÓGICA DEL ATLAS
% ===============================================================

\section{Construcción Metodológica del Atlas Solar Territorial}

\figuraopcional{figuras_informe/figura_01_pipeline_atlas_solar.png}
{Esquema metodológico general y cadena computacional del Atlas Solar de la Provincia de Córdoba: desde la ingesta de datos satelitales GOES hasta los productos territoriales de potencial fotovoltaico.}
{fig:pipeline_general}

\subsection{Dominio espacial y grilla territorial}

El dominio de trabajo corresponde al territorio de la Provincia de Córdoba. La representación espacial se apoya en la grilla regular definida desde las etapas iniciales del proyecto, con separación aproximada de $0.018^\circ$, equivalente nominalmente a unos 2~km sobre el dominio provincial.

El buffer externo utilizado durante el procesamiento funciona únicamente como margen técnico para evitar pérdidas de información en bordes durante operaciones de reproyección, interpolación y muestreo satelital. Los productos provinciales finales se recortan al polígono administrativo de Córdoba.

La cantidad final de celdas interiores utilizada en los productos territoriales es:

\[
N_{\text{celdas}} = 48\,645.
\]

La grilla extendida utilizada durante el procesamiento contiene 59\,837 posiciones canónicas; de ellas, 48\,645 se encuentran dentro del territorio provincial y son las utilizadas para las estadísticas oficiales.

Es importante distinguir la resolución de la grilla del Atlas de la resolución nativa del sensor ABI \cite{schmit2017}. La información satelital visible se extrae de las escenas GOES y se proyecta/interpola sobre la grilla canónica del Atlas; por lo tanto, la expresión ``grilla de 2 km'' se refiere al soporte territorial del producto y no a la resolución nativa del canal visible del satélite.

\subsection{Base satelital y período de inferencia}

Para cada celda y cada instante operativo se utilizaron observaciones del canal visible C02 del sensor ABI a bordo de GOES-16/19, junto con las variables geométricas y radiométricas requeridas por el modelo satelital.

La cadena territorial utiliza:

\begin{itemize}
    \item reflectancia visible observada por ABI;
    \item brillo de cielo despejado $B_0$;
    \item término nuboso relativo $F_r$;
    \item geometría solar;
    \item irradiancia de cielo despejado $GHI_{cs}$;
    \item variables temporales requeridas por el componente HGB.
\end{itemize}

Todas las fuentes se expresan en UTC y se procesan bajo una convención temporal homogénea. El detalle de calibración y selección de la arquitectura predictiva fue desarrollado en el Tercer Informe; aquí se utiliza directamente la configuración consolidada para la inferencia territorial.

\subsection{Regionalización solar}

La relación entre reflectancia satelital, nubosidad e irradiancia no es necesariamente homogénea en toda la Provincia de Córdoba. Para representar esta heterogeneidad se utiliza una regionalización operativa basada en la red terrestre calibrada.

La clasificación parte de indicadores agregados por estación. Las variables empleadas son:

\begin{itemize}
    \item latitud;
    \item longitud;
    \item altitud;
    \item GHI diurna media calibrada;
    \item índice de claridad medio $K_t$.
\end{itemize}

Antes del clustering, las variables se estandarizan para impedir que las diferencias de escala numérica condicionen artificialmente las distancias. El agrupamiento se realiza mediante K-Means y el número de grupos se evalúa mediante Silhouette Score.

Sobre la base terrestre consolidada, la configuración adoptada corresponde a:

\[
k=3.
\]

Las tres regiones operativas se interpretan como:

\begin{enumerate}[label=R\arabic*:]
    \setcounter{enumi}{-1}
    \item Sierras Altas;
    \item Sierras Bajas / zona de transición;
    \item Llanura Pampeana.
\end{enumerate}

Estas denominaciones facilitan la interpretación física del agrupamiento, pero las regiones no deben interpretarse como una zonificación climática oficial o definitiva. Su función es reducir la heterogeneidad interna del entrenamiento y organizar la inferencia territorial del modelo satelital.

\figuraopcional[0.72\textwidth]{figuras_informe/figura_02_regionalizacion_red_estaciones.png}
{Regionalización solar operativa de la red de estaciones de superficie (k=3) adoptada para la inferencia territorial.}
{fig:regionalizacion}

\figuraopcional{figuras_informe/figura_03_caracterizacion_regiones_solares.png}
{Caracterización biofísica y radiométrica de las tres regiones solares operativas ($k=3$): comparación de altitud topográfica, irradiancia global diurna media y factor de calibración de la red de estaciones terrestres.}
{fig:caracterizacion_regiones}

\subsection{Asignación espacial de regiones a la grilla}

El \textit{clustering} se define originalmente sobre la red territorial que corresponde a estimadores puntuales. Para utilizar los modelos regionales sobre toda la provincia es necesario propagar esa clasificación a cada celda territorial.

La asignación se realiza mediante un clasificador de tres vecinos más próximos (3-NN) en coordenadas métricas planas (POSGAR 2007 / Argentina 3 --- EPSG:5345). Para evitar distorsiones asociadas al uso directo de latitud y longitud, las distancias espaciales se calculan en el sistema proyectado oficial utilizado por el proyecto para las operaciones métricas provinciales.

El clasificador utiliza ponderación por distancia. Para una celda $\mathbf{p}$ se obtienen pesos regionales:

\[
w_r(\mathbf{p}), \qquad r\in\{0,1,2\},
\]

con

\[
\sum_{r=0}^{2} w_r(\mathbf{p}) = 1.
\]

Estos pesos deben interpretarse como un consenso espacial del vecindario 3-NN; no representan probabilidades climáticas calibradas ni una estimación formal de incertidumbre.

\figuraopcional{figuras_informe/figura_04_propagacion_regiones_grid_3nn.png}
{Propagación territorial de las regiones solares a la grilla nominal de 2 km mediante clasificador 3-NN ponderado por distancia: (a) región dominante discreta y (b) grado de consenso espacial vecinal.}
{fig:propagacion_regiones}

\subsection{Asignación Hard y combinación Blended}

Se evaluaron dos estrategias para aplicar los modelos regionales sobre la grilla.

En la configuración \textit{Hard}, cada celda queda asociada exclusivamente a la región dominante:

\[
r^*(\mathbf{p})=\arg\max_r w_r(\mathbf{p}),
\]

y la estimación utiliza únicamente el modelo correspondiente a $r^*$.

En la configuración \textit{Blended}, las predicciones de los tres modelos regionales se combinan utilizando los pesos del clasificador:

\[
\widehat{K_t}(\mathbf{p},t)
=
\sum_{r=0}^{2}
w_r(\mathbf{p})\,
\widehat{K_{t,r}}(\mathbf{p},t).
\]

La estrategia Hard tiene la ventaja de una interpretación regional discreta, pero puede introducir cambios artificialmente abruptos en los límites entre clases. La estrategia Blended conserva la especialización regional y, al mismo tiempo, permite una transición continua entre áreas vecinas.

La comparación cuantitativa entre ambas configuraciones se presenta en la Tabla~\ref{tab:hard_vs_blended}.

\begin{table}[htbp]
\centering
\small
\caption{Sensibilidad espacial y comparación cuantitativa: Modo Hard vs. Modo Blended.}
\label{tab:hard_vs_blended}
\begin{tabular}{lcc}
\toprule
Métrica & Hard & Blended \\
\midrule
GHI medio provincial (kWh/m$^2$/año) & 1885.9 & 1888.3 \\
Desvío espacial de GHI (kWh/m$^2$/año) & 80.2 & 78.5 \\
$Y_f$ fijo provincial (kWh/kWp/año) & 1643.4 & 1645.2 \\
Diferencia relativa media provincial & \multicolumn{2}{c}{+0.13\%} \\
Diferencia media absoluta espacial & \multicolumn{2}{c}{0.41\% (rango local: $-2.26\%$ a $+2.18\%$)} \\
\bottomrule
\end{tabular}
\end{table}

\figuraopcional{figuras_informe/figura_05_hard_vs_blended.png}
{Evaluación de sensibilidad espacial: (a) asignación discreta Hard, (b) asignación continua Blended y (c) diferencia porcentual relativa entre ambos esquemas, evidenciando que las discrepancias se concentran en las zonas de transición sin perturbar el balance energético provincial.}
{fig:hard_vs_blended}

La configuración operativa adoptada para el Atlas es el modo Blended. El modo Hard se conserva como análisis de sensibilidad.

\subsection{Arquitectura de inferencia territorial}

El modelo operativo conserva la estructura híbrida seleccionada en el Tercer Informe. Para cada región solar y estación del año se dispone de una parametrización específica. Con tres regiones y cuatro estaciones se obtiene una matriz de:

\[
N_{\text{modelos}}=3\times4=12
\]

configuraciones regional-estacionales.

Para cada celda y tiempo, el flujo de inferencia es:

\begin{enumerate}[label=\arabic*)]
    \item calcular la geometría solar;
    \item calcular $GHI_{cs}$ mediante la formulación operativa de cielo despejado;
    \item obtener reflectancia ABI y construir las variables $B_0$ y $F_r$;
    \item aplicar el modelo BD-JPT correspondiente;
    \item calcular la corrección residual mediante HGB;
    \item obtener una predicción regional de $K_t$;
    \item combinar las tres predicciones regionales mediante los pesos Blended;
    \item limitar el resultado al intervalo físico-operativo adoptado;
    \item reconstruir GHI.
\end{enumerate}

Un factor determinante en la precisión del modelo satelital corresponde a la sincronización temporal entre la observación geoestacionaria instantánea y la integral horaria registrada por los piranómetros terrestres. Los ensayos de desarrollo demostraron que la asignación instantánea en extremo de intervalo (Estrategia A) introducía un desacople dinámico respecto de la respuesta física de la atmósfera. En su lugar, el Atlas adoptó la \textbf{Estrategia Temporal C}, en la cual la reflectancia satelital se sincroniza como el promedio centrado entre los extremos del paso horario $[t, t+1\text{ h})$:

\[
\overline{R}_t = \frac{R_t + R_{t+1}}{2},
\]

garantizando una correspondencia radiométrica óptima con la irradiancia acumulada en superficie.

\figuraopcional{figuras_informe/figura_07_alineacion_temporal_A_vs_C.png}
{Diagnóstico de sincronización temporal entre la observación satelital y la red de superficie: comparación entre la alineación instantánea en extremo de intervalo (Estrategia A) y el promedio centrado de paso horario (Estrategia C).}
{fig:alineacion_temporal}

La predicción híbrida regional puede expresarse como:

\[
K_{t,r}^{\mathrm{hyb}}
=
\texttt{clip}
\left[
K_{t,r}^{\mathrm{BDJPT}}+
\Delta K_{t,r}^{\mathrm{HGB}},
0,\;1.15
\right].
\]

Luego,

\[
K_t^{\mathrm{Atlas}}
=
\sum_{r=0}^{2}w_rK_{t,r}^{\mathrm{hyb}},
\]

y finalmente,

\[
GHI_{\mathrm{Atlas}}
=
K_t^{\mathrm{Atlas}}\,GHI_{cs}.
\]

El modelo de cielo despejado utilizado en esta etapa cumple un rol de normalización física del problema satelital y no reemplaza el uso de REST2 como referencia radiométrica en la calibración terrestre.

\subsection{Agregación temporal y estadísticas multianuales}

La salida instantánea u horaria del modelo se convierte en energía integrada utilizando la duración efectiva de cada intervalo $\Delta t_i$.

Para una celda $\mathbf{p}$ y un año $y$:

\[
H_{\mathrm{GHI},y}(\mathbf{p})
=
\frac{1}{1000}
\sum_{i\in y}
GHI_i(\mathbf{p})\,\Delta t_i,
\]

donde $GHI_i$ se expresa en W/m$^2$, $\Delta t_i$ en horas y el resultado en kWh/m$^2$/año.

El campo multianual final 2020--2025 se calcula como el promedio de los seis totales anuales:

\[
\overline{H}_{\mathrm{GHI}}^{2020-2025}(\mathbf{p})
=
\frac{1}{6}
\sum_{y=2020}^{2025}
H_{\mathrm{GHI},y}(\mathbf{p}).
\]

Esta operación evita que un año con mayor cantidad de registros disponibles reciba automáticamente un peso mayor que los demás.

Para caracterizar el ciclo anual se utilizan medias diarias mensuales. Para un mes $m$:

\[
\overline{H}_{m,\text{día}}(\mathbf{p})
=
\frac{\sum_{d\in m} H_d(\mathbf{p})}{N_{d,m}},
\]

expresada en kWh/m$^2$/día.

Esta métrica permite comparar meses de diferente duración sin confundir el efecto del número de días con la intensidad radiativa.

\subsection{Horas Sol Pico}

Las Horas Sol Pico (HSP) representan la energía solar diaria expresada como el número equivalente de horas durante las cuales la irradiancia es constante e igual a 1~kW/m$^2$.

Por lo tanto, una irradiación diaria de:

\[
5~\text{kWh/m}^2\text{/día}
\]

equivale numéricamente a:

\[
5~\text{HSP/día}.
\]

Para cada celda, el producto HSP se calcula a partir de la irradiación diaria media del período 2020--2025.

\subsection{Control de cobertura y consistencia}

La agregación temporal debe distinguir entre irradiancia físicamente nula y ausencia de información. Las horas nocturnas no constituyen datos faltantes: su contribución energética es nula por definición. En cambio, una escena satelital ausente durante un intervalo diurno representa una pérdida real de cobertura.

Para cada estadístico se registra su cobertura efectiva, entendida como la fracción de intervalos diurnos teóricamente disponibles que cuentan con una estimación válida y participan en el cálculo.

Se integraron todos los días del período analizado:

\[
\boxed{2192/2192\ \text{días}}
\]

Dentro del dominio operacional del modelo, definido por $GHI_{cs}>20\ \mathrm{W/m^2}$ y cenit solar aparente inferior a $87^\circ$, la disponibilidad de intervalos diurnos válidos fue superior al 98\%. Los huecos aislados de hasta una hora fueron completados mediante interpolación lineal, mientras que se exigió una completitud mínima del 90\% para productos anuales y del 85\% para productos mensuales. Las celdas o años que no alcanzaron este umbral fueron excluidas del cálculo multianual para evitar transformar problemas de disponibilidad en sesgos artificiales de irradiación anual acumulada.

\figuraopcional[0.72\textwidth]{figuras_informe/figura_17_cobertura_efectiva.png}
{Cobertura efectiva de datos diurnos válidos en la grilla territorial de 2 km.}
{fig:cobertura}

\subsection{Validación y controles numéricos}

La validación del modelo satelital fue desarrollada en detalle en el Tercer Informe de avances. Para la versión territorial consolidada se conservan controles adicionales orientados a asegurar que la aplicación masiva del modelo no introduzca inconsistencias respecto de la configuración validada.

Las métricas consolidadas para la estrategia temporal C se resumen a continuación. La validación temporal corresponde al test independiente enero--agosto de 2021, mientras que Pilar CMP21 se reporta como referencia radiométrica puntual:

\begin{table}[htbp]
\centering
\small
\caption{Métricas consolidadas de validación para la arquitectura satelital híbrida (Estrategia Temporal C).}
\label{tab:validacion_modelo}
\begin{tabular}{lccc}
\toprule
Conjunto & RMSE GHI (W/m$^2$) & MAE GHI (W/m$^2$) & $R^2$ GHI \\
\midrule
CV espacial (train 2020) & 115.96 & 78.45 & 0.8400 \\
Test temporal OOS 2021 & 120.68 & 77.63 & 0.7939 \\
Pilar CMP21 & 75.17 & 55.90 & 0.9166 \\
\bottomrule
\end{tabular}
\end{table}

Todos los errores de GHI se expresan en W/m$^2$. Estas métricas corresponden a la cadena corregida y no deben compararse directamente con métricas intermedias de informes anteriores obtenidas sobre configuraciones de datos diferentes.

\figuraopcional{figuras_informe/figura_06_validacion_modelo_final.png}
{Validación numérica consolidada del modelo satelital híbrido (Estrategia Temporal C): (a) validación cruzada espacial (CV 2020) sobre la red terrestre ($170\,254$ pares), (b) test independiente fuera de muestra temporal (OOS 2021, $198\,609$ pares) y (c) validación de alta precisión frente a la estación radiométrica patrón Pilar CMP21 ($5\,122$ pares horarios).}
{fig:validacion_final}


% ===============================================================
% 3. RECURSO SOLAR TERRITORIAL
% ===============================================================

\section{Distribución Territorial del Recurso Solar}

\subsection{Irradiación Global Horizontal}

El producto principal del Atlas es la irradiación global horizontal anual media del período 2020--2025:

\[
\overline{H}_{\mathrm{GHI}}
\quad [\text{kWh/m}^2\text{/año}].
\]

Este mapa representa la energía solar total recibida sobre una superficie horizontal durante un año medio construido a partir de los seis años disponibles. Valores mayores indican una mayor disponibilidad global de energía solar, integrando tanto la componente directa proyectada sobre el plano horizontal como la radiación difusa.

La media provincial para el período 2020--2025 es:

\[
\boxed{
\overline{H}_{\mathrm{GHI,prov}}^{2020-2025}
=
1889.3
\;\text{kWh/m}^2\text{/año}
}
\]

El rango espacial provincial correspondiente es:

\[
1081.9
\leq
\overline{H}_{\mathrm{GHI}}^{2020-2025}
\leq
2054.7
\;\text{kWh/m}^2\text{/año}.
\]

Este promedio equivale a aproximadamente:

\[
5.17\;\text{HSP/día},
\]

con un rango espacial de 2.96 a 5.63 HSP/día.

\subsection{Estimación derivada de DNI y DHI}

El modelo satelital principal estima GHI. La irradiancia directa normal (DNI) y la irradiancia difusa horizontal (DHI) no son variables observadas independientemente en la red provincial utilizada para el Atlas. Se obtienen como productos derivados mediante descomposición radiativa.

Por esta razón, los mapas de DNI y DHI deben interpretarse como estimaciones consistentes con el campo de GHI, la geometría solar y el modelo de descomposición adoptado, pero no como observaciones directas ni como productos validados de manera independiente.

Los valores provinciales consolidados se expresan a continuación:

\[
\boxed{
\overline{H}_{\mathrm{DNI,prov}}^{2020-2025}
=
1985.2
\;\text{kWh/m}^2\text{/año}
}
\]

con un rango territorial de 396.6 a 2279.2 kWh/m$^2$/año, mientras que:

\[
\boxed{
\overline{H}_{\mathrm{DHI,prov}}^{2020-2025}
=
604.9
\;\text{kWh/m}^2\text{/año}
}
\]

con un rango territorial de 564.5 a 910.2 kWh/m$^2$/año.

\figuraopcional{figuras_informe/figura_08_recurso_solar_ghi_dni_dhi.png}
{Distribución territorial de los componentes principales del recurso solar en Córdoba: (a) Irradiación Global Horizontal (GHI), (b) Irradiación Directa Normal (DNI) y (c) Irradiación Difusa Horizontal (DHI).}
{fig:recurso_solar}

\subsection{Ciclo mensual del recurso}

Para describir el ciclo estacional se calculan medias diarias mensuales en kWh/m$^2$/día utilizando los años del período 2020--2025.

Este formato permite interpretar directamente cuánta energía solar recibe, en promedio, una superficie horizontal durante un día típico de cada mes.

\figuraopcional{figuras_informe/figura_09_ciclo_mensual_anual_completo.png}
{Ciclo mensual completo del recurso solar en la Provincia de Córdoba: los 12 meses representados en irradiación diaria media (kWh/m$^2$/día) bajo una escala radiométrica homogénea y compartida.}
{fig:ciclo_mensual}

En la Tabla~\ref{tab:ghi} se muestran los valores GHI y HSP/día medios para toda la provincia:

\begin{table}[htbp]
\centering
\small
\caption{Ciclo mensual medio provincial de irradiación global horizontal (GHI) y Horas Sol Pico (HSP).}
\label{tab:ghi}
\begin{tabular}{lcc}
\toprule
Mes & GHI media diaria (kWh/m$^2$/día) & HSP/día \\
\midrule
Enero & 7.35 & 7.35 \\
Febrero & 6.71 & 6.71 \\
Marzo & 5.25 & 5.25 \\
Abril & 4.25 & 4.25 \\
Mayo & 3.27 & 3.27 \\
Junio & 2.76 & 2.76 \\
Julio & 3.02 & 3.02 \\
Agosto & 4.15 & 4.15 \\
Septiembre & 5.17 & 5.17 \\
Octubre & 6.14 & 6.14 \\
Noviembre & 6.77 & 6.77 \\
Diciembre & 7.30 & 7.30 \\
\midrule
\textbf{Media Anual} & \textbf{5.17} & \textbf{5.17} \\
\bottomrule
\end{tabular}
\end{table}

\figuraopcional[0.85\textwidth]{figuras_informe/figura_10_perfil_mensual_provincial.png}
{Perfil estacional de la irradiación global horizontal media diaria en la Provincia de Córdoba, con banda de dispersión espacial percentil 10 a percentil 90 y desglose por región operativa.}
{fig:perfil_mensual}

\subsection{Variabilidad interanual observada}

El período 2020--2025 permite cuantificar la estabilidad temporal observada del campo a partir de seis años completos. Este intervalo permite evaluar la variabilidad interanual interna del producto, aunque sigue siendo insuficiente para definir una normal climatológica de largo plazo.

Para cada celda se calcula el coeficiente de variación interanual, definido como:

\[
CV(\mathbf{p})
=
100
\frac{
\sigma_y\left[H_{\mathrm{GHI},y}(\mathbf{p})\right]
}{
\overline{H}_{\mathrm{GHI}}(\mathbf{p})
}.
\]

El $CV$ expresa, en porcentaje, cuánto varían los totales anuales alrededor de su media multianual. Valores bajos indican que los años analizados presentan totales similares; valores más altos identifican sectores donde la variabilidad interanual observada es mayor.

Para el período analizado, la media provincial de $CV$ es:

\[
\boxed{
CV_{\mathrm{prov}}=2.14\%
}
\]

con valores territoriales comprendidos entre 0.47\% y 16.99\%.

\figuraopcional[0.72\textwidth]{figuras_informe/figura_14_cv_interanual.png}
{Coeficiente de variación interanual ($CV$, \%) de la irradiación solar anual en la Provincia de Córdoba.}
{fig:cv_interanual}

La sensibilidad del campo multianual a la incorporación sucesiva de nuevos años se mantuvo por debajo del 1\% en todas las etapas evaluadas:

\[
|\Delta_2|=0.98\%,\qquad
|\Delta_3|=0.57\%,\qquad
|\Delta_4|=0.96\%,\qquad
|\Delta_5|=0.46\%,\qquad
|\Delta_6|=0.15\%,
\]

con una media provincial de $\Delta_6=+0.07\%$ (rango territorial entre $-1.56\%$ y $+6.41\%$).

La secuencia no se interpreta como una convergencia monótona, dado que la magnitud del cambio puede aumentar al incorporar un año meteorológicamente diferente. Sin embargo, los cambios observados en las etapas sucesivas evaluadas permanecen por debajo del 1\%, lo que indica una elevada estabilidad del promedio multianual del período 2020--2025.

\figuraopcional{figuras_informe/figura_13_sensibilidad_multianual.png}
{Sensibilidad del promedio multianual a la incorporación sucesiva de nuevos años ($\Delta_2$, $\Delta_3$, $\Delta_4$, $\Delta_5$ y $\Delta_6$), evidenciando la estabilidad y convergencia del campo acumulado.}
{fig:sensibilidad_multianual}

Las correlaciones espaciales interanuales fueron elevadas, con valores comprendidos entre $r=0.881$ y $r=0.950$, indicando una persistencia importante de la organización territorial del recurso solar entre años.

\figuraopcional{figuras_informe/figura_15_estabilidad_espacial_anios.png}
{Comparación espacial año por año de la irradiación global horizontal anual, demostrando la persistencia del gradiente territorial a lo largo del período.}
{fig:anios_comparacion}

\figuraopcional[0.65\textwidth]{figuras_informe/figura_16_correlacion_espacial_interanual.png}
{Matriz de correlación espacial interanual de Pearson entre los mapas anuales de irradiación global horizontal.}
{fig:matriz_correlacion}

\subsection{Patrones espaciales y diferencias regionales}

El análisis territorial se complementa con estadísticas agregadas por las tres regiones solares operativas. Para cada región se calculan media, desviación espacial y percentiles de GHI, DNI, DHI y HSP.

\begin{table}[htbp]
\centering
\footnotesize
\setlength{\tabcolsep}{3.5pt}
\caption{Distribución del recurso solar y potencial fotovoltaico por región operativa (período 2020--2025).}
\label{tab:recurso_regional}
\begin{tabular}{lrrrrrrrrr}
\toprule
Región & Celdas & GHI & DNI & DHI & $Y_f$ fijo & $Y_f$ trk & CF fijo & CF trk & Ganancia \\
 & & \multicolumn{3}{c}{(kWh/m$^2$/año)} & \multicolumn{2}{c}{(kWh/kWp/año)} & (\%) & (\%) & (\%) \\
\midrule
R0 Sierras Altas & 5\,586 & 1898.6 & 2030.8 & 595.4 & 1659.9 & 1999.6 & 18.95\% & 22.83\% & 20.45\% \\
R1 Sierras Bajas / Transición & 17\,677 & 1922.9 & 2033.1 & 607.8 & 1671.2 & 2020.1 & 19.08\% & 23.06\% & 20.85\% \\
R2 Llanura Pampeana & 25\,382 & 1863.8 & 1941.9 & 604.9 & 1622.9 & 1952.0 & 18.53\% & 22.28\% & 20.27\% \\
\midrule
\textbf{Provincia} & \textbf{48\,645} & \textbf{1889.3} & \textbf{1985.2} & \textbf{604.9} & \textbf{1644.7} & \textbf{1982.2} & \textbf{18.77\%} & \textbf{22.63\%} & \textbf{20.51\%} \\
\bottomrule
\end{tabular}
\end{table}

Se observa que la Región 1 (Sierras Bajas y Transición) presenta el mayor promedio de radiación solar global (1\,922.9~kWh/m$^2$/año) debido a su elevada heliofanía relativa y moderada nubosidad convectiva, mientras que la Llanura Pampeana (R2) presenta un promedio ligeramente menor (1\,863.8~kWh/m$^2$/año) con una dispersión espacial muy baja (desviación estándar de solo 34.1~kWh/m$^2$/año).

\subsection{Singularidades espaciales y efecto del albedo superficial}

El uso principal de un canal visible implica una limitación conocida en superficies excepcionalmente brillantes. La señal satelital utiliza la reflectancia como proxy de nubosidad: en condiciones normales, un incremento marcado de reflectancia respecto del fondo despejado suele asociarse con presencia de nubes. Sin embargo, superficies de albedo muy alto pueden generar una respuesta similar aun bajo cielo despejado.

En Córdoba, el caso más relevante corresponde a las Salinas Grandes. La elevada reflectancia de la costra salina puede superar ampliamente el valor modelado de brillo de fondo provincial $B_0$, aumentando artificialmente el término nuboso relativo $F_r = \max(R^2 - B_0^2, 0)$ bajo cielo completamente despejado y produciendo una subestimación local de $K_t$ y GHI.

Este comportamiento no debe interpretarse como una característica real del recurso solar sin una revisión específica. A continuación se detallan las estadísticas del área problemática:

\begin{itemize}
    \item cantidad de celdas identificadas en la zona anómala: 359 celdas, equivalentes al 0.74\% de las celdas provinciales del Atlas ($359 / 48\,645 = 0.738\%$);
    \item sesgo aparente respecto del entorno: $-585.0\text{ kWh/m}^2/\text{año}$ ($-30.2\%$, con valores medios en el salar de $1\,350\text{ kWh/m}^2/\text{año}$ frente a $1\,935\text{ kWh/m}^2/\text{año}$ en la zona circundante a misma latitud, y mínimos puntuales de $1\,081.9\text{ kWh/m}^2/\text{año}$);
    \item causa física demostrada: el albedo superficial de la costra salina seca ($R \approx 0.40$--$0.50$ bajo cielo despejado) sobrepasa la envolvente provincial homogénea de cielo claro ($B_0 \approx 0.20$), lo que produce un contraste nuboso aparente positivo $F_r > 0$ que induce al modelo a inferir nubosidad inexistente;
    \item tratamiento adoptado: mantener la anomalía identificada y documentada como limitación localizada del esquema visible monocanal; evaluar una máscara específica de superficies de alto albedo en versiones posteriores.
\end{itemize}

\figuraopcional{figuras_informe/figura_18_diagnostico_salinas_grandes.png}
{Diagnóstico de la singularidad espacial asociada al alto albedo de Salinas Grandes: (a) localización geográfica del salar en el noroeste provincial, (b) detalle cartográfico de la depresión radiométrica y (c) transecta zonal oeste--este a latitud $30.05^\circ$S cuantificando el sesgo de $-30.2\%$ confinado a la cuenca salina.}
{fig:salinas}


% ===============================================================
% 4. MODELADO Y POTENCIAL FOTOVOLTAICO TERRITORIAL
% ===============================================================

\section{Modelado y Potencial Fotovoltaico Territorial}

\subsection{Alcance de la estimación fotovoltaica}

El mapa de GHI describe el recurso solar disponible sobre una superficie horizontal, pero no representa directamente la energía que produciría una instalación fotovoltaica.

Para estimar producción eléctrica potencial se requiere una cadena adicional de modelos:

\[
GHI
\rightarrow
(DNI,DHI)
\rightarrow
G_{\mathrm{POA}}
\rightarrow
T_{\mathrm{cell}}
\rightarrow
P_{\mathrm{DC}}
\rightarrow
P_{\mathrm{AC}}
\rightarrow
E_{\mathrm{AC}}.
\]

Cada transformación agrega supuestos. En consecuencia, los mapas de producción fotovoltaica se interpretan como escenarios tecnológicos normalizados y no como una predicción específica para una planta real.

\subsection{Descomposición GHI $\rightarrow$ DNI + DHI: modelo de Erbs}

La relación geométrica entre las componentes radiativas es:

\[
GHI = DNI\cos\theta_z + DHI.
\]

A partir de GHI solamente, esta ecuación no permite determinar de manera única DNI y DHI. Para cerrar el problema se utiliza el modelo empírico de Erbs, que estima la fracción difusa a partir del índice de claridad radiométrico relativo a la irradiancia extraterrestre.

Debe distinguirse este índice del $K_t$ utilizado en el modelo satelital híbrido. En la etapa satelital se utiliza:

\[
K_t^{cs}=\frac{GHI}{GHI_{cs}},
\]

mientras que la descomposición de Erbs utiliza un índice basado en irradiancia extraterrestre horizontal:

\[
k_t^{Erbs}
=
\frac{GHI}{I_{0,h}}.
\]

A partir de $k_t^{Erbs}$ se estima la fracción difusa $k_d$, y luego:

\[
DHI = k_d\,GHI,
\]

\[
DNI=
\frac{GHI-DHI}{\cos\theta_z},
\]

aplicando controles físicos para sol bajo, irradiancias pequeñas y resultados negativos o no plausibles.

\subsection{Transposición al plano del generador: modelo de Perez}

Una vez obtenidas DNI y DHI, la irradiancia sobre el plano del generador se calcula mediante el modelo anisotrópico de Perez.

La irradiancia total en el plano puede expresarse como:

\[
G_{\mathrm{POA}}
=
G_{b,\mathrm{POA}}
+
G_{d,\mathrm{sky}}
+
G_{d,\mathrm{ground}}.
\]

El término directo depende del ángulo de incidencia entre los rayos solares y la normal al módulo. El término difuso de cielo es resuelto mediante la formulación de Perez, que representa la anisotropía del cielo, mientras que la componente reflejada incorpora la contribución del suelo.

Se analizan dos configuraciones de referencia:

\subsubsection{Sistema fijo}

El sistema fijo se define mediante:

\begin{itemize}
    \item inclinación: $30^\circ$;
    \item orientación: norte geográfico;
    \item potencia nominal normalizada: 1~kWp.
\end{itemize}

Esta configuración permite construir un escenario homogéneo para toda la provincia y comparar el efecto territorial del recurso solar sin introducir optimizaciones locales de inclinación.

\subsubsection{Seguidor horizontal de un eje}

La configuración con seguimiento utiliza:

\begin{itemize}
    \item eje horizontal norte--sur;
    \item rotación efectiva este--oeste;
    \item ángulo máximo: $60^\circ$;
    \item backtracking activado;
    \item Ground Coverage Ratio (GCR): 0.35;
    \item potencia nominal normalizada: 1~kWp.
\end{itemize}

El backtracking reduce pérdidas geométricas por sombreado mutuo entre filas durante horas de sol bajo.

\subsection{Modelo térmico NOCT}

La eficiencia eléctrica del módulo disminuye cuando aumenta la temperatura de celda. Para representar este efecto se utiliza un modelo térmico NOCT simplificado:

\[
T_{\mathrm{cell}}
=
T_{\mathrm{amb}}
+
\left(
\frac{NOCT-20}{800}
\right)
G_{\mathrm{POA}},
\]

con:

\[
NOCT = 45^\circ\mathrm{C}.
\]

En la implementación territorial, $T_{\mathrm{amb}}$ se parametriza mediante temperaturas estacionales representativas de Córdoba:

\begin{center}
\begin{tabular}{lc}
\toprule
Estación & $T_{\mathrm{amb}}$ \\
\midrule
Verano (DJF) & $24^\circ$C \\
Otoño (MAM) & $18^\circ$C \\
Invierno (JJA) & $11^\circ$C \\
Primavera (SON) & $19^\circ$C \\
\bottomrule
\end{tabular}
\end{center}

Esta simplificación permite incorporar el efecto térmico de primer orden, pero no representa explícitamente variaciones horarias de temperatura ambiente, viento, configuración de montaje o inercia térmica del módulo.

\subsection{Conversión fotovoltaica mediante PVWatts}

La potencia de corriente continua se estima mediante la formulación PVWatts:

\[
P_{\mathrm{DC}}
=
P_{\mathrm{DC,0}}
\left(
\frac{G_{\mathrm{POA}}}{1000}
\right)
\left[
1+\gamma_P(T_{\mathrm{cell}}-25)
\right].
\]

La simulación se normaliza a una potencia nominal de:

\[
P_{\mathrm{DC,0}}=1~\mathrm{kWp}.
\]

El coeficiente térmico adoptado es:

\[
\gamma_P=-0.0035~^\circ\mathrm{C}^{-1}.
\]

La conversión DC/AC utiliza el modelo de inversor PVWatts con eficiencia nominal:

\[
\eta_{\mathrm{inv,nom}}=0.96.
\]

Adicionalmente se aplica una eficiencia global de balance de sistema:

\[
\eta_{\mathrm{BOS}}=0.90,
\]

equivalente a pérdidas adicionales globales del 10\% para representar de manera simplificada cableado, \textit{mismatch}, suciedad, disponibilidad y otras pérdidas no resueltas explícitamente.

Los supuestos anteriores definen un escenario tecnológico común para todas las celdas y permiten comparar el potencial relativo del territorio. No sustituyen un estudio de ingeniería de una instalación específica.

\subsection{Rendimiento final}

El rendimiento final $Y_f$ representa la energía eléctrica anual entregada por unidad de potencia pico instalada:

\[
Y_f
=
\frac{E_{\mathrm{AC}}}{P_{\mathrm{nom}}},
\]

con unidades:

\[
[\;Y_f\;]
=
\text{kWh/kWp/año}.
\]

Al normalizar la simulación a 1~kWp, $Y_f$ coincide numéricamente con la energía anual AC producida por el sistema de referencia.

Los valores provinciales definitivos se resumen a continuación:

\[
\boxed{
Y_{f,\mathrm{fijo}}^{2020-2025}
=
1644.7
\;\text{kWh/kWp/año}
}
\]

con un rango territorial de 870.5 a 1792.6 kWh/kWp/año, y:

\[
\boxed{
Y_{f,\mathrm{tracker}}^{2020-2025}
=
1982.2
\;\text{kWh/kWp/año}
}
\]

con un rango territorial de 995.6 a 2167.5 kWh/kWp/año.

\subsection{Factor de capacidad}

El factor de capacidad expresa qué fracción de la energía máxima teórica produciría el sistema si operara a potencia nominal durante todas las horas del año:

\[
CF
=
100
\frac{Y_f}{8760}.
\]

Para mantener una convención uniforme y directamente comparable entre celdas y años, los factores de capacidad reportados por el Atlas se normalizan respecto de 8760 horas por año. Esta convención coincide con la definición utilizada por la cadena de cálculo de los productos consolidados; la pequeña diferencia introducida en años bisiestos no modifica la interpretación territorial del indicador.

Los valores provinciales consolidados para 2020--2025 son:

\[
\boxed{
CF_{\mathrm{fijo}}=18.77\%
}
\]

con un rango territorial de 9.94\% a 20.46\%, y:

\[
\boxed{
CF_{\mathrm{tracker}}=22.63\%
}
\]

con un rango territorial de 11.37\% a 24.74\%.

\subsection{Mapas de irradiación en plano y rendimiento eléctrico}

Para el período 2020--2025, la irradiación anual media sobre el plano fijo de $30^\circ$ norte es:

\[
\boxed{
POA_{\mathrm{fijo}}=2016.6\;\text{kWh/m}^2\text{/año}
}
\]

con un rango territorial de 1042.3 a 2206.0 kWh/m$^2$/año. Para el seguidor horizontal de un eje se obtiene:

\[
\boxed{
POA_{\mathrm{tracker}}=2461.6\;\text{kWh/m}^2\text{/año}
}
\]

con un rango territorial de 1201.5 a 2701.6 kWh/m$^2$/año.

Los productos principales de esta sección son:

\begin{itemize}
    \item POA anual medio para sistema fijo;
    \item POA anual medio para tracker;
    \item $Y_f$ fijo;
    \item $Y_f$ tracker;
    \item factor de capacidad fijo;
    \item factor de capacidad tracker.
\end{itemize}

\figuraopcional{figuras_informe/figura_11_poa_y_rendimiento_fotovoltaico.png}
{Irradiación sobre el plano del generador y rendimiento fotovoltaico específico para las configuraciones fija y con seguimiento: (a) POA fijo, (b) POA tracker, (c) $Y_f$ fijo y (d) $Y_f$ tracker.}
{fig:potencial_fv}


% ===============================================================
% 5. GANANCIA POR SEGUIMIENTO Y ANÁLISIS REGIONAL
% ===============================================================

\section{Ganancia por Seguimiento y Análisis Regional}

\subsection{Definición de ganancia por seguimiento}

La ventaja energética del seguidor de un eje se expresa mediante la ganancia relativa respecto del sistema fijo:

\[
G_{\mathrm{tracker}}
=
100
\left(
\frac{Y_{f,\mathrm{tracker}}}
{Y_{f,\mathrm{fijo}}}
-1
\right).
\]

Esta métrica permite comparar directamente ambas configuraciones independientemente del nivel absoluto de producción de cada celda.

La media provincial final es:

\[
\boxed{
\overline{G}_{\mathrm{tracker}}^{2020-2025}
=
20.51\%
}
\]

El rango territorial de la ganancia es de 14.38\% a 23.36\%.

\figuraopcional{figuras_informe/figura_12_cf_y_ganancia_tracker.png}
{Factor de capacidad para el sistema fijo y el seguidor de un eje, junto con la ganancia porcentual del seguimiento respecto del sistema fijo de $30^\circ$ norte.}
{fig:ganancia_tracker}

\subsection{Interpretación física de la ganancia}

La ganancia del seguimiento no depende únicamente de la cantidad total de GHI. Un tracker modifica de manera continua la orientación del plano del generador y aprovecha especialmente la componente directa de la irradiancia durante una fracción mayor del día.

Por este motivo, la ganancia espacial debe analizarse conjuntamente con:

\begin{itemize}
    \item DNI estimada;
    \item fracción difusa;
    \item geometría solar;
    \item distribución temporal de la nubosidad;
    \item disponibilidad radiativa durante mañana y tarde.
\end{itemize}

Dos regiones con valores anuales similares de GHI pueden presentar ganancias de seguimiento diferentes si difiere la proporción de irradiancia directa y difusa.

\subsection{Comparación entre regiones solares}

Las tres regiones operativas permiten sintetizar el comportamiento territorial del recurso y del potencial fotovoltaico. Los valores correspondientes se resumen en la Tabla~\ref{tab:ghi_region}.

\begin{table}[htbp]
\centering
\small
\caption{Comparación de indicadores del recurso solar y potencial fotovoltaico por región solar operativa.}
\label{tab:ghi_region}
\begin{tabular}{lccccccc}
\toprule
Región & GHI & DNI & $Y_f$ fijo & $Y_f$ trk. & CF fijo & CF trk. & Ganancia \\
\midrule
R0 & 1898.6 & 2030.8 & 1659.9 & 1999.6 & 18.95\% & 22.83\% & 20.45\% \\
R1 & 1922.9 & 2033.1 & 1671.2 & 2020.1 & 19.08\% & 23.06\% & 20.85\% \\
R2 & 1863.8 & 1941.9 & 1622.9 & 1952.0 & 18.53\% & 22.28\% & 20.27\% \\
\midrule
\textbf{Provincia} & \textbf{1889.3} & \textbf{1985.2} & \textbf{1644.7} & \textbf{1982.2} & \textbf{18.77\%} & \textbf{22.63\%} & \textbf{20.51\%} \\
\bottomrule
\end{tabular}
\end{table}

La interpretación regional revela que la Región 1 (Sierras Bajas y Transición) combina la mayor irradiación directa con la mayor ganancia media por seguimiento ($+20.85\%$), convirtiéndola en la zona con mayor productividad anual tanto en sistemas fijos ($1\,671.2$~kWh/kWp/año) como con seguimiento ($2\,020.1$~kWh/kWp/año).

\subsection{Comportamiento estacional de la ganancia}

Además del mapa anual, la ganancia puede calcularse por mes o estación del año:

\[
G_{\mathrm{tracker},m}
=
100
\left(
\frac{Y_{f,\mathrm{tracker},m}}
{Y_{f,\mathrm{fijo},m}}
-1
\right).
\]

El desglose estacional provincial y regional de la ganancia por seguimiento se detalla en la Tabla~\ref{tab:ganancia_estacional}.

\begin{table}[htbp]
\centering
\small
\caption{Comportamiento estacional de la ganancia relativa por seguimiento (\%).}
\label{tab:ganancia_estacional}
\begin{tabular}{lcccc}
\toprule
Estación & Provincial & R0 (Sierras Altas) & R1 (Transición) & R2 (Llanura) \\
\midrule
Verano (DJF) & 24.8\% & 24.5\% & 25.1\% & 24.6\% \\
Primavera (SON) & 21.5\% & 21.4\% & 21.8\% & 21.3\% \\
Otoño (MAM) & 17.6\% & 17.7\% & 17.9\% & 17.4\% \\
Invierno (JJA) & 14.2\% & 14.4\% & 14.6\% & 14.0\% \\
\midrule
\textbf{Anual} & \textbf{20.5\%} & \textbf{20.5\%} & \textbf{20.9\%} & \textbf{20.3\%} \\
\bottomrule
\end{tabular}
\end{table}

El máximo beneficio del seguimiento ocurre en verano ($24.8\%$), cuando la prolongada duración diurna y las elevadas alturas solares permiten a los seguidores capturar de manera óptima la irradiancia directa temprana y tardía mediante la rotación este--oeste.


% ===============================================================
% 6. PRODUCTOS CARTOGRÁFICOS Y GEOESPACIALES
% ===============================================================

\section{Productos Cartográficos y Base Geoespacial}

\subsection{Productos principales}

El Atlas se organiza en productos de recurso solar, productos fotovoltaicos y productos diagnósticos.

\begin{longtable}{p{0.31\textwidth}p{0.20\textwidth}p{0.39\textwidth}}
\toprule
Producto & Unidad & Interpretación \\
\midrule
\endfirsthead
\toprule
Producto & Unidad & Interpretación \\
\midrule
\endhead

GHI anual medio &
kWh/m$^2$/año &
Energía solar anual recibida sobre una superficie horizontal. Producto base del Atlas. \\

DNI anual media &
kWh/m$^2$/año &
Componente directa normal estimada a partir de GHI mediante descomposición radiativa. \\

DHI anual media &
kWh/m$^2$/año &
Componente difusa horizontal estimada a partir de GHI. \\

HSP &
HSP/día &
Número equivalente de horas diarias a 1~kW/m$^2$. Numéricamente equivalente a la irradiación diaria en kWh/m$^2$/día. \\

POA fijo $30^\circ$ N &
kWh/m$^2$/año &
Energía radiativa anual incidente sobre el plano de un sistema fijo de referencia. \\

POA tracker &
kWh/m$^2$/año &
Energía radiativa anual incidente sobre el plano móvil del seguidor de un eje. \\

$Y_f$ fijo &
kWh/kWp/año &
Energía eléctrica anual normalizada por kWp instalado para el sistema fijo. \\

$Y_f$ tracker &
kWh/kWp/año &
Energía eléctrica anual normalizada por kWp instalado para el sistema con seguimiento. \\

Factor de capacidad fijo &
\% &
Fracción equivalente del año durante la cual el sistema fijo produciría a potencia nominal. \\

Factor de capacidad tracker &
\% &
Fracción equivalente del año durante la cual el tracker produciría a potencia nominal. \\

Ganancia tracker/fijo &
\% &
Incremento relativo de rendimiento anual del tracker respecto del sistema fijo. \\

Variabilidad interanual &
\% &
Coeficiente de variación de los seis totales anuales del período 2020--2025. \\

Regionalización solar &
categoría / pesos &
Región operativa dominante y pesos de consenso espacial utilizados por la inferencia Blended. \\

Cobertura efectiva &
\% &
Fracción de intervalos válidos que contribuyen a cada estadístico agregado. \\

\bottomrule
\end{longtable}

\subsection{Formatos geoespaciales y especificaciones técnicas}

Los mapas ráster finales se exportan en formatos geoespaciales compatibles con herramientas GIS y servicios web. El formato principal es Cloud-Optimized GeoTIFF (COG), acompañado por tablas vectoriales GeoPackage y Parquet.

Las especificaciones técnicas y la estructura de archivos entregables se detallan en el Apéndice~\ref{sec:especificaciones}.

\subsection{Visualización interactiva}

Los archivos ráster COG permiten publicar los productos sin requerir que el usuario final utilice QGIS. Se diseñó un visor web interactivo basado en Leaflet y WebGL con capacidad de consulta de valores puntuales por celda, perfiles y metadatos.

\figuraopcional{figuras_informe/figura_19_visor_atlas_solar.png}
{Plataforma web interactiva del Atlas Solar de Córdoba con soporte para capas ráster COG y consulta territorial por celda.}
{fig:visor}

\subsection{Reproducibilidad}

Todos los productos territoriales se generan mediante \textit{scripts} reproducibles a partir de los datasets consolidados. La separación entre datos de entrada, modelos, productos intermedios y mapas finales permite repetir la generación completa del Atlas cuando se incorporen nuevos años, nuevas calibraciones o versiones actualizadas de los modelos.


% ===============================================================
% 7. CONCLUSIONES
% ===============================================================

\section{Conclusiones}

\subsection{Distribución provincial del recurso solar}

Para el período 2020--2025, la irradiación global horizontal media provincial es de 1889.3 kWh/m$^2$/año, equivalente a 5.17 HSP/día. La DNI derivada alcanza 1985.2 kWh/m$^2$/año y la DHI 604.9 kWh/m$^2$/año. El GHI presenta un rango territorial de 1081.9 a 2054.7 kWh/m$^2$/año.

La variabilidad interanual observada es reducida a escala provincial: el coeficiente de variación medio para 2020--2025 es 2.14\% (con rango entre 0.47\% y 16.99\%). La sensibilidad del campo medio a la incorporación de años sucesivos se mantiene por debajo del 1\% en las etapas evaluadas, con una variación al incorporar 2025 de apenas $|\Delta_6|=0.15\%$ (media provincial $+0.07\%$). Las correlaciones espaciales entre $r=0.881$ y $r=0.950$ entre los mapas anuales indican que la organización territorial del recurso es más persistente que las variaciones interanuales de su magnitud absoluta.

\subsection{Potencial de generación fotovoltaica}

La irradiación media sobre el plano de un sistema fijo de $30^\circ$ norte alcanza 2016.6 kWh/m$^2$/año, mientras que el seguidor horizontal de un eje alcanza 2461.6 kWh/m$^2$/año. Bajo los supuestos tecnológicos adoptados, esto se traduce en rendimientos específicos de 1644.7 kWh/kWp/año para el sistema fijo y 1982.2 kWh/kWp/año para el tracker.

Los factores de capacidad medios correspondientes son 18.77\% y 22.63\%, respectivamente. Estos indicadores representan escenarios tecnológicos normalizados y no sustituyen estudios de ingeniería específicos para una planta o emplazamiento determinado.

\subsection{Beneficio territorial del seguimiento en un eje}

La ganancia energética media del seguidor respecto del sistema fijo es 20.51\% para el período 2020--2025. A escala territorial, la ganancia varía entre 14.38\% y 23.36\%, mostrando que el beneficio del seguimiento no es espacialmente uniforme y depende no solo del GHI total sino también de la composición directa/difusa y de la distribución temporal del recurso.

\subsection{Principales contrastes regionales}

La regionalización operativa divide las 48\,645 celdas provinciales en 5\,586 celdas de Sierras Altas (11.5\%), 17\,677 de Sierras Bajas / Transición (36.3\%) y 25\,382 de Llanura Pampeana (52.2\%). La Región 1 destaca como la más productiva de la provincia, alcanzando rendimientos de 2\,020.1~kWh/kWp/año con seguimiento y 1\,671.2~kWh/kWp/año en configuración fija.

La principal singularidad espacial identificada corresponde a Salinas Grandes, donde 359 celdas, equivalentes al 0.74\% de las celdas provinciales del Atlas, presentan un sesgo negativo de $-30.2\%$ respecto del entorno circundante, debido a la confusión radiométrica entre la costra salina de alto albedo y nubosidad en el esquema visible monocanal. Esta limitación se mantiene explícitamente documentada en el producto.

\subsection{Productos generados para el Atlas Solar}

La cadena territorial produce mapas de GHI, DNI, DHI, HSP, irradiación en plano fijo y con seguimiento, rendimiento específico, factor de capacidad, ganancia de seguimiento, regionalización solar y variabilidad interanual. Los productos científicos se exportan en formatos geoespaciales interoperables (COG, GPKG) y se acompañan de figuras cartográficas y un visor interactivo para facilitar su consulta.


% ===============================================================
% ANEXOS
% ===============================================================

\appendix

\section{Ecuaciones y convenciones de referencia}

\subsection{Índice de claridad del modelo satelital}

\[
K_t^{cs}=\frac{GHI}{GHI_{cs}}.
\]

\subsection{Reconstrucción híbrida}

\[
K_t^{\mathrm{hyb}}
=
\operatorname{clip}
\left(
K_t^{\mathrm{BDJPT}}+\Delta K_t^{\mathrm{HGB}},
0,\;1.15
\right).
\]

\subsection{Descomposición radiativa}

\[
GHI=DNI\cos\theta_z+DHI.
\]

\subsection{Irradiancia en plano}

\[
G_{\mathrm{POA}}
=
G_{b,\mathrm{POA}}
+
G_{d,\mathrm{sky}}
+
G_{d,\mathrm{ground}}.
\]

\subsection{Temperatura de celda}

\[
T_{\mathrm{cell}}
=
T_{\mathrm{amb}}
+
\left(
\frac{NOCT-20}{800}
\right)
G_{\mathrm{POA}}.
\]

\subsection{Potencia DC}

\[
P_{\mathrm{DC}}
=
P_{\mathrm{DC,0}}
\left(
\frac{G_{\mathrm{POA}}}{1000}
\right)
\left[
1+\gamma_P(T_{\mathrm{cell}}-25)
\right].
\]

\subsection{Rendimiento final}

\[
Y_f=\frac{E_{\mathrm{AC}}}{P_{\mathrm{nom}}}.
\]

\subsection{Factor de capacidad}

\[
CF=100\frac{Y_f}{8760}.
\]

\subsection{Ganancia del seguimiento}

\[
G_{\mathrm{tracker}}
=
100
\left(
\frac{Y_{f,\mathrm{tracker}}}{Y_{f,\mathrm{fijo}}}
-1
\right).
\]

% ============================================================
\section{Tabla maestra de resultados consolidados}
% ============================================================

\begin{longtable}{p{0.50\textwidth}p{0.42\textwidth}}
\toprule
Resultado & Valor final \\
\midrule
\endfirsthead
\toprule
Resultado & Valor final \\
\midrule
\endhead

Número de celdas provinciales & 48\,645 (59\,837 posiciones en la grilla extendida) \\
Cobertura temporal efectiva & 2192/2192 días analizados (disponibilidad diurna $>98\%$ dentro del dominio operacional) \\
GHI provincial medio & 1889.3 kWh/m$^2$/año (5.17 HSP/día) \\
DNI provincial media & 1985.2 kWh/m$^2$/año \\
DHI provincial media & 604.9 kWh/m$^2$/año \\
HSP provincial medio & 5.17 HSP/día \\
POA fijo provincial & 2016.6 kWh/m$^2$/año \\
POA tracker provincial & 2461.6 kWh/m$^2$/año \\
$Y_f$ fijo provincial & 1644.7 kWh/kWp/año \\
$Y_f$ tracker provincial & 1982.2 kWh/kWp/año \\
CF fijo provincial & 18.77\% \\
CF tracker provincial & 22.63\% \\
Ganancia tracker/fijo provincial & 20.51\% (rango 14.38--23.36\%) \\
CV interanual provincial & 2.14\% (rango 0.47--16.99\%) \\
Sensibilidad temporal sucesiva & $|\Delta_2|=0.98\%$, $|\Delta_3|=0.57\%$, $|\Delta_4|=0.96\%$, $|\Delta_5|=0.46\%$, $|\Delta_6|=0.15\%$ (media $+0.07\%$) \\
Correlaciones espaciales interanuales & $r=0.881$ a $r=0.950$ (mín. $r_{2021,2023}=0.881$; máx. $r_{2023,2024}=0.950$) \\
Rango territorial GHI provincial & 1081.9 a 2054.7 kWh/m$^2$/año \\
Rango territorial $Y_f$ & Fijo: 870.5 a 1792.6; Tracker: 995.6 a 2167.5 kWh/kWp/año \\
Celdas afectadas por singularidad Salinas Grandes & 359 (0.74\% de las celdas provinciales del Atlas) \\
Sesgo de irradiación en Salinas Grandes & $-585.0$ kWh/m$^2$/año ($-30.2\%$ respecto al entorno regional) \\
Sensibilidad Hard vs. Blended (GHI) & Hard: 1885.9 kWh/m$^2$/a vs. Blended: 1888.3 kWh/m$^2$/a (+0.13\%) \\
Sensibilidad Hard vs. Blended ($Y_f$) & Hard: 1643.4 kWh/kWp/a vs. Blended: 1645.2 kWh/kWp/a (+0.11\%) \\
Dispersión espacial Hard vs. Blended & $\sigma_{\mathrm{Hard}}=80.2$ kWh/m$^2$/a vs. $\sigma_{\mathrm{Blended}}=78.5$ kWh/m$^2$/a \\
Métricas de validación (CV Espacial train 2020) & RMSE 115.96 W/m$^2$, MAE 78.45 W/m$^2$, $R^2=0.8400$ \\
Métricas de validación (Test OOS 2021) & RMSE 120.68 W/m$^2$, MAE 77.63 W/m$^2$, $R^2=0.7939$ \\
Métricas de validación (Pilar CMP21) & RMSE 75.17 W/m$^2$, MAE 55.90 W/m$^2$, $R^2=0.9166$ \\
Estadísticas R0 (Sierras Altas) & 5\,586 celdas (11.5\%); GHI 1898.6, DNI 2030.8, DHI 595.4, $Y_f$ 1659.9 / 1999.6, Gan. 20.45\% \\
Estadísticas R1 (Sierras Bajas / Trans.) & 17\,677 celdas (36.3\%); GHI 1922.9, DNI 2033.1, DHI 607.8, $Y_f$ 1671.2 / 2020.1, Gan. 20.85\% \\
Estadísticas R2 (Llanura Pampeana) & 25\,382 celdas (52.2\%); GHI 1863.8, DNI 1941.9, DHI 604.9, $Y_f$ 1622.9 / 1952.0, Gan. 20.27\% \\
\bottomrule
\end{longtable}

% ============================================================
\section{Diagnóstico de evolución histórica del modelo satelital}
\label{sec:anexo_evolucion}
% ============================================================

A título de documentación técnica interna sobre las etapas de desarrollo previas a la consolidación del modelo final, la Figura~\ref{fig:anexo_evolucion} conserva una de las evaluaciones intermedias de corrección residual sobre la estación Pilar, previa a la adopción definitiva de la Estrategia C y del ensamble regional-estacional consolidado.

\figuraopcional[0.85\textwidth]{figuras_informe/anexo/anexo_07_diagnostico_evolucion_modelo.png}
{Diagnóstico de evolución y corrección del modelo satelital preliminar evaluado en la estación Pilar (etapa intermedia de desarrollo técnico).}
{fig:anexo_evolucion}

% ============================================================
\section{Especificaciones de productos geoespaciales}
\label{sec:especificaciones}
% ============================================================

Todos los productos ráster generados por la cadena del Atlas Solar se distribuyen en formato Cloud-Optimized GeoTIFF (COG), estructurados con teselas internas (\textit{tiles}) de $128 \times 128$, compresión DEFLATE sin pérdida y niveles de reducción de resolución (\textit{overviews}) 2, 4 y 8.

Para facilitar la localización y transferencia a usuarios externos (organismos públicos, evaluadores energéticos, desarrolladores de proyectos y la comunidad técnica y académica), los archivos ráster se entregan organizados en carpetas temáticas normalizadas con sus respectivos catálogos explicativos (\path{README.md}):

{\small
\begin{longtable}{p{0.44\textwidth}p{0.16\textwidth}p{0.14\textwidth}p{0.22\textwidth}}
\toprule
Ruta y Archivo ráster (COG) & Variable & Unidad & Descripción \\
\midrule
\endfirsthead
\toprule
Ruta y Archivo ráster (COG) & Variable & Unidad & Descripción \\
\midrule
\endhead

\textbf{01\_recurso\_solar\_multianual/} & & & \\
\path{cordoba_ghi_anual_medio_2020_2025_2km_cog.tif} & Irradiación GHI & kWh/m$^2$/año & GHI medio sexenal 2020--2025 \\
\path{cordoba_dni_anual_medio_2020_2025_2km_cog.tif} & Irradiación DNI & kWh/m$^2$/año & DNI derivada mediante Erbs \\
\path{cordoba_dhi_anual_medio_2020_2025_2km_cog.tif} & Irradiación DHI & kWh/m$^2$/año & DHI derivada mediante Erbs \\
\path{cordoba_psh_anual_medio_2020_2025_2km_cog.tif} & Horas Sol Pico & HSP/día & Promedio diario anual equivalente \\
\midrule
\textbf{02\_potencial\_fotovoltaico\_multianual/} & & & \\
\path{cordoba_poa_fijo_anual_medio_2020_2025_2km_cog.tif} & POA Fijo 30°N & kWh/m$^2$/año & Irradiación plano fijo 30° N \\
\path{cordoba_poa_tracker_anual_medio_2020_2025_2km_cog.tif} & POA Tracker 1-eje & kWh/m$^2$/año & Irradiación seguidor 1-eje N-S \\
\path{cordoba_yf_fijo_anual_medio_2020_2025_2km_cog.tif} & Rendimiento $Y_f$ Fijo & kWh/kWp/año & Producción específica en plano fijo \\
\path{cordoba_yf_tracker_anual_medio_2020_2025_2km_cog.tif} & Rendimiento $Y_f$ Tracker & kWh/kWp/año & Producción específica con seguidor \\
\path{cordoba_cf_fijo_anual_medio_2020_2025_2km_cog.tif} & Factor de Capacidad & \% & CF sistema fijo (base 8760 h) \\
\path{cordoba_cf_tracker_anual_medio_2020_2025_2km_cog.tif} & Factor de Capacidad & \% & CF sistema seguidor (base 8760 h) \\
\path{cordoba_ganancia_tracker_anual_medio_2020_2025_2km_cog.tif} & Ganancia Seguidor & \% & Incremento relativo seguidor vs fijo \\
\midrule
\textbf{03\_ciclo\_mensual\_multianual/} & & & \\
\path{cordoba_ghi_mensual_01_a_12_2020_2025_2km_cog.tif} & GHI Mensual (12 capas) & kWh/m$^2$/mes & Irradiación acumulada por mes \\
\midrule
\textbf{04\_series\_anuales\_individuales/} & & & \\
\path{{2020..2025}/cordoba_*_anual_{año}_2km_cog.tif} & Variables anuales & Diversas & Series individuales completas (6 años) \\
\midrule
\textbf{05\_variabilidad\_y\_convergencia/} & & & \\
\path{cordoba_cv_interanual_2020_2025_2km_cog.tif} & Coef. de Variación & \% & Variabilidad interanual sexenal \\
\path{cordoba_delta_2_a_6_pct_2km_cog.tif} & Sensibilidades $\Delta_k$ & \% & Estabilización asintótica temporal \\
\midrule
\textbf{06\_regionalizacion\_y\_zonas\_homogeneas/} & & & \\
\path{cordoba_regiones_solares_k3_2km_cog.tif} & Regiones k=3 & Categoría & Zonas solares homogéneas (0, 1, 2) \\
\bottomrule
\end{longtable}
}

Parámetros comunes de los archivos ráster:
\begin{itemize}
    \item \textbf{Sistema de Coordenadas Geográficas (CRS):} EPSG:4326 (WGS 84), con soporte nativo de proyección métrica POSGAR 2007 / Argentina 3 (EPSG:5345);
    \item \textbf{Resolución espacial nominal:} $0.018^\circ \times 0.018^\circ$ ($\approx 2.0$~km);
    \item \textbf{Dimensiones de la matriz provincial:} 328 filas $\times$ 244 columnas ($48\,645$ celdas válidas interiores);
    \item \textbf{Valor de No-Dato (\textit{NoData}):} $-9999.0$;
    \item \textbf{Método de agregación temporal:} Integración horaria trapezoidal/rectangular diaria $\rightarrow$ agregación mensual $\rightarrow$ promedio multianual ponderado uniforme;
    \item \textbf{Versión del producto:} Release 5B.3.0\_DEFINITIVO.
\end{itemize}


\begin{thebibliography}{99}

\bibitem{gueymard2008}
Gueymard, C. A. (2008).
REST2: High-performance solar radiation model for cloudless-sky irradiance, illuminance, and photosynthetically active radiation: Validation with a benchmark dataset.
\textit{Solar Energy}, 82(3), 272--285.

\bibitem{tarpley1979}
Tarpley, J. D. (1979).
Estimating incident solar radiation at the surface from geostationary satellite data.
\textit{Journal of Applied Meteorology}, 18(9), 1172--1181.

\bibitem{alonso2012}
Alonso-Suárez, R., Abal, G., Siri, R., \& Musé, P. (2012).
Brightness-dependent Tarpley model for global solar radiation estimation using GOES satellite images: Application to Uruguay.
\textit{Solar Energy}, 86(11), 3205--3215.

\bibitem{ineichen2002}
Ineichen, P., \& Perez, R. (2002).
A new airmass independent formulation for the Linke turbidity coefficient.
\textit{Solar Energy}, 73(3), 151--157.

\bibitem{holmgren2018}
Holmgren, W. F., Hansen, C. W., \& Mikofski, M. A. (2018).
pvlib python: A Python package for modeling solar energy systems.
\textit{Journal of Open Source Software}, 3(29), 884.

\bibitem{erbs1982}
Erbs, D. G., Klein, S. A., \& Duffie, J. A. (1982).
Estimation of the diffuse radiation fraction for hourly, daily and monthly-average global radiation.
\textit{Solar Energy}, 28(4), 293--302.

\bibitem{perez1990}
Perez, R., Ineichen, P., Seals, R., Michalsky, J., \& Stewart, R. (1990).
Modeling daylight availability and irradiance components from direct and global irradiance.
\textit{Solar Energy}, 44(5), 271--289.

\bibitem{dobos2014}
Dobos, A. P. (2014).
\textit{PVWatts Version 5 Manual}.
National Renewable Energy Laboratory, NREL/TP-6A20-62641.

\bibitem{faiman2008}
Faiman, D. (2008).
Assessing the outdoor operating temperature of photovoltaic modules.
\textit{Progress in Photovoltaics: Research and Applications}, 16(4), 307--315.

\bibitem{schmit2017}
Schmit, T. J., Griffith, P., Gunshor, M. M., Daniels, J. M., Goodman, S. J., \& Lebair, W. J. (2017).
A closer look at the ABI on the GOES-R series.
\textit{Bulletin of the American Meteorological Society}, 98(4), 681--698.

\end{thebibliography}

\end{document}
